\documentclass[11pt]{article}

\usepackage[margin=1in]{geometry}
\usepackage{amsmath,amssymb,amsthm}
\usepackage{booktabs}
\usepackage{microtype}
\usepackage[hidelinks]{hyperref}

\newtheorem{proposition}{Proposition}
\newtheorem{conjecture}{Conjecture}
\newcommand{\E}{\mathbb{E}}
\newcommand{\Prb}{\mathop{\mathrm{Pr}}}

\title{Why the Current Fixed-Lane ORAM Stash Depends on $N$}
\author{}
\date{July 2026}

\begin{document}
\maketitle

\begin{abstract}
Circuit ORAM has an $N$-uniform, exponentially decaying snapshot stash tail
in its proved parameter regime.  The current fixed-lane ORAM transition does
not exhibit this behavior.  For a binary tree of depth
$L=\log_2 N$, experiments and a product-state flow calculation indicate that
one saturated lane carries only $\Theta(L^{-1/2})$ useful load.  Consequently,
maintaining a fixed operating margin requires approximately
$Z=\Omega(\sqrt{\log_2 N})$ lanes.  This note gives exact conservation
identities, a mean-field explanation of the square-root law, empirical
validation, and a comparison with Circuit ORAM.  The conservation identities
are exact; the stationary square-root law is a strongly supported conjecture,
not yet a security theorem.
\end{abstract}

\section{Setting}

Consider the current binary fixed-lane transition with $N=2^L$ logical
blocks and $Z$ physical lanes per bucket.  Every access reads the path given
by the requested block's old uniformly random label, removes the block,
assigns it a fresh independent label, inserts it into the stash, admits stash
blocks into free root lanes, and performs one lane-local eviction on the old
path.

The physical paths are IID uniform under a fixed repeated-permutation logical
schedule.  The observed $N$-dependence is therefore not caused by biased or
insufficiently random path selection.  It is caused by the amount of useful
flow that the fixed-lane transition can sustain through a depth-$L$ routing
network.

\section{Exact conservation laws}

Let $S_t$ be post-writeback stash occupancy after operation $t$.  During the
next operation, let $I_t$ equal one if the requested block is removed from the
tree and zero if it was already in the stash.  Let $A_t$ be the number of
stash blocks admitted into free root lanes.  Block conservation gives the
following identity without any probabilistic assumptions.

\begin{proposition}[Root conservation]
\label{prop:root}
For every operation,
\[
  S_{t+1}-S_t=I_t-A_t.
\]
In stationarity,
\[
  \E[A_t]=\Prb[I_t=1].
\]
\end{proposition}

There is an analogous identity across every tree cut.  Fix a depth $d$, and
let $U_{d,t}$ be the number of blocks stored at depths $0,\ldots,d$.  Let
$R_{d,t}$ indicate removal of the requested block from this upper region, and
let $C_{d,t}$ count blocks moved downward across the cut below depth $d$.
Lane eviction never moves blocks upward, so
\[
  U_{d,t+1}-U_{d,t}=A_t-R_{d,t}-C_{d,t}.
\]
Combining stationary balance with Proposition~\ref{prop:root} yields
\begin{proposition}[Cut conservation]
\label{prop:cut}
In stationarity,
\[
  \E[C_{d,t}]
  =\Prb[\text{the requested block is in the tree below depth }d].
\]
\end{proposition}

The upper $d$ levels have only $Z(2^{d+1}-1)$ slots.  Hence a small-stash
ORAM containing $N$ blocks must place almost all blocks below every fixed
shallow cut as $N$ grows.  Proposition~\ref{prop:cut} then requires nearly
one aggregate downward crossing per operation through almost every shallow
cut.  Nominal tree capacity by itself is insufficient; the eviction process
must provide the required flow at every depth.

\section{Why lower occupancy does not create first-order slack}

Suppose a small-stash system divides one unit of aggregate flow among $Z$
lanes.  Each lane must carry approximately
\[
  q\simeq \frac{1}{Z}
\]
blocks per operation.  At binary depth $d$, a block sees the correct next-bit
eviction path once per $2^{d+1}$ global operations, while arrivals are split
among $2^d$ nodes.  In dilute flow, Little's-law reasoning therefore gives a
typical lane-slot occupancy
\[
  \rho\simeq 2q.
\]

Thus adding lanes really does lower occupancy.  However, a useful move also
requires an occupied source.  The ideal binary flow from one lane is
$\rho/2\simeq q$, exactly the offered load.  The first-order benefit of more
holes is cancelled by the first-order loss of occupied sources, leaving no
constant service margin.

For an independent Bernoulli path state with occupied probability $\rho$, a
direct expansion of the production mask's local cut current gives
\begin{equation}
\label{eq:current}
  J_2(\rho)
  =\frac{\rho}{2}-\frac{3}{64}\rho^3+O(\rho^4).
\end{equation}
The quadratic term cancels because the mask can jump over one incompatible
occupied blocker.  Substituting $\rho\simeq2q$ gives
\begin{equation}
\label{eq:flow-q}
  J_2(2q)=q-\frac{3}{8}q^3+O(q^4).
\end{equation}
Consequently, one routing level loses a relative $O(q^2)$ fraction of ideal
flow.

A lane storing a non-negligible fraction of $N$ blocks must route blocks down
to depth $L-O(\log L)$, so the loss is accumulated over $\Theta(L)$ prefix
stages.  Equations~\eqref{eq:current}--\eqref{eq:flow-q} therefore suggest
that the relevant dimensionless congestion parameter is
\[
  Lq^2.
\]
Nonvanishing end-to-end flow requires $Lq^2=O(1)$, giving
\begin{equation}
\label{eq:sqrt-law}
  q_L=\Theta(L^{-1/2}).
\end{equation}
This is a mean-field scaling argument rather than a proof for the correlated
stationary ORAM process.

An equivalent intuition views a vacancy created deep in a lane as prefix
credit that must propagate back to the root.  Random left/right demand and
eviction opportunities produce an approximately zero-drift reflected walk.
Its typical imbalance or maximum over $L$ stages is $\Theta(\sqrt L)$, so the
usable load is its reciprocal.  Fixed lane identity prevents imbalances and
vacancies from being pooled across lanes.

\section{Measured per-lane law}

Production-equivalent saturated one-lane experiments used direct operating
system randomness, random initial labels, fresh independent labels on every
access, and at least $64N$ warm-up operations.  Define
\[
  q_L=\frac{\E[\text{blocks resident in one lane tree}]}{N}.
\]
The measured products were
\begin{center}
\begin{tabular}{c@{\qquad}c}
\toprule
$L$ & $q_L\sqrt L$ \\
\midrule
10 & 0.6317 \\
12 & 0.6221 \\
14 & 0.6222 \\
16 & 0.6246 \\
18 & 0.6317 \\
20 & 0.6327 \\
\bottomrule
\end{tabular}
\end{center}
These results support
\begin{equation}
\label{eq:fitted-q}
  q_L\simeq\frac{c}{\sqrt L},\qquad c\simeq0.627.
\end{equation}
The fitted constant is strikingly close to
$\sqrt{\pi/8}=0.626657\ldots$.  A random-walk maximum heuristic can produce
this value, but the required mapping and its factor of two have not been
proved.  The constant must therefore be regarded as empirical.

When the global stash is backlogged, the $Z$ marginal lanes behave
approximately as parallel saturated copies.  This gives the fluid
approximation
\begin{equation}
\label{eq:stash-law}
  \boxed{
  \frac{\E[S]}{N}
  \simeq
  \left(1-\frac{0.627Z}{\sqrt{\log_2N}}\right)_+.}
\end{equation}
The corresponding crossover height is
\begin{equation}
\label{eq:crossover}
  L_{\mathrm{crit}}(Z)\simeq(0.627Z)^2,
  \qquad
  Z_{\mathrm{crit}}(L)\simeq1.595\sqrt L.
\end{equation}

\begin{center}
\begin{tabular}{ccc}
\toprule
$Z$ & Predicted $L_{\mathrm{crit}}$ & Observed behavior \\
\midrule
5 & 9.8  & degradation near $N=2^{10}$ \\
6 & 14.2 & degradation near $N=2^{14}$ \\
7 & 19.3 & small through $2^{18}$, large at $2^{20}$ \\
8 & 25.2 & small but increasingly heavy-tailed through $2^{20}$ \\
9 & 31.8 & small through $2^{20}$; no asymptotic guarantee \\
\bottomrule
\end{tabular}
\end{center}

The abrupt transition is ordinary queue stability.  For example, at $Z=7$
the fitted aggregate capacity is $1.034$ at $L=18$ but $0.981$ at $L=20$.
A few percent change moves the queue from slightly negative to slightly
positive drift.  Since $N=2^L$, adding one lane can postpone the transition
by a very large multiplicative factor in $N$.

\section{Direct $Z=6$ and $Z=7$ results}

The current block-conserving simulator produced the following mean
post-writeback stash occupancies.
\begin{center}
\begin{tabular}{r@{\qquad}r@{\qquad}r}
\toprule
$N$ & $Z=6$ & $Z=7$ \\
\midrule
$2^{10}$ & 0.773 & 0.076 \\
$2^{12}$ & 6.288 & 0.222 \\
$2^{14}$ & 207.214 & 0.762 \\
$2^{16}$ & 4314.083 & 3.668 \\
$2^{18}$ & 29660.358 & 32--41 \\
$2^{20}$ & --- & 9416.649 \\
\bottomrule
\end{tabular}
\end{center}
At $N=2^{20}$, the $Z=7$ measured range was 1542--16682 after a
$2^{26}$-operation warm-up.  Thus its excellent small-$N$ behavior was a
finite-height effect, not an $N$-uniform stash bound.

At $N=2^{20}$, current $Z=8$ had mean stash 0.873 but transient insertion
demand exceeded 20 with probability approximately $1.10\%$.  Current $Z=9$
had mean 0.049 and exceedance probability approximately
$1.33\times10^{-4}$.  These finite-run probabilities are snapshot estimates,
not lifetime failure bounds or cryptographic extrapolations.

\section{Greedy fixed-lane control}

A corrected, conserving deepest-hole transition improves the per-lane
constant but does not change the observed exponent:
\[
  q_L^{\mathrm{greedy}}\simeq\frac{0.82}{\sqrt L}.
\]
At $N=2^{20}$, greedy $Z=6$ and $Z=7$ had mean post-stash occupancies 2.529
and 0.095, respectively, compared with 9416.649 for current $Z=7$.
The predicted greedy crossover heights are approximately 24 for $Z=6$ and
33 for $Z=7$.  Greedy movement therefore matters substantially, but the
fixed-lane routing evidence still suggests $Z=\Omega(\sqrt{\log_2N})$ for
binary trees.

\section{Why Circuit ORAM is different}

Circuit ORAM's proved snapshot stash bound has the form
\[
  \Prb[S>R]\le C\alpha^R,
  \qquad \alpha<1,
\]
with constants independent of $N$ in the theorem's parameter regime.
Therefore its expected snapshot stash is $O(1)$ independently of $N$; only a
separate lifetime union bound forces $R$ to grow with the number of operations
or desired security parameter.

Circuit ORAM does not route vacancies through permanently separated channels.
It pools bucket slots, chooses the deepest compatible candidate from the
stash and path, moves directly into deep vacancies, and performs two eviction
paths per logical access.  These mechanisms provide an $N$-uniform service
margin.  Lane ORAM does have $Z$ potential parallel chains, but the useful
capacity of each fixed chain decreases with routing depth because small local
losses accumulate at every prefix stage.

\section{General branching factors}

Binary branching is unusually favorable.  In the same independent path-state
model, general branching factor $B$ gives
\[
  J_B(\rho)
  =\frac{\rho}{B}
   -\frac{(B-1)(B-2)}{B^3}\rho^2
   +O(\rho^3).
\]
The quadratic loss cancels only for $B=2$.  Since dilute occupancy satisfies
$\rho\simeq Bq$, for $B>2$ the relative loss per level is $O(q)$ rather than
$O(q^2)$.  Across $L=\log_BN$ levels this suggests
\[
  q_L=\Theta(1/L),
  \qquad Z=\Omega(\log_BN).
\]
Existing $B=4$ data fit $q_LL\simeq0.363$, and independent $Z$ values infer
the same marginal per-lane capacity.  This explains why the tested $B=4$ and
$B=8$ configurations perform worse than their raw number of simultaneous
lane moves suggests.

\section{Status of the analysis}

Propositions~\ref{prop:root} and~\ref{prop:cut} are exact consequences of
block conservation.  Equation~\eqref{eq:current} is exact for an IID
Bernoulli product-state model of a loaded path.  The extension from that local
model to the correlated stationary ORAM process, the square-root stationary
law, and the numerical constant in Equation~\eqref{eq:fitted-q} remain
conjectural.  Their unusually accurate agreement across tree sizes and lane
counts makes them a strong explanation of the observed behavior, but not a
replacement for a formal stash-security proof.

\end{document}
